\section{Sonuçlar}
\label{sec:results}

\subsection{Dondurulmuş eleme ve transfer kazanımı}
Tüm tekli/ikili/üçlü yapılandırmalar önce 0.~katta dondurulmuş rejimde
sıralanmıştır. Dondurulmuş rejimde en güçlü \emph{tekli} omurga Swin-T
(makro-F1 $0{,}576$), en güçlü \emph{üçlü} ise ağırlıklı füzyon ($0{,}565$)
olmuştur; tek başına BEiT-B en zayıf omurgadır ($0{,}522$). Son blokların ince
ayarı, eleme adayı her yapılandırmada makro-F1'i artırmıştır
(Tablo~\ref{tab:transfer}); en büyük kazanım ikili ve üçlü füzyon
yapılandırmalarında görülmüştür. Bu, küçük veri rejiminde bile son dönüştürücü
bloklarının uç (endoskopi) alanına uyarlanmasının değerli olduğunu gösterir.
Tablo~\ref{tab:frozen} tüm yapılandırmaların dondurulmuş sıralamasını verir.

\begin{table}[H]
  \centering
  \caption{0.~kat dondurulmuş eleme sıralaması (makro-F1'e göre).
  V=ViT-B, S=Swin-T, B=BEiT-B.}
  \label{tab:frozen}
  \begin{tabular}{cllcc}
    \toprule
    \# & Omurga & Füzyon & Doğruluk & Makro-F1 \\
    \midrule
    1  & S     & yok      & $0{,}8704$ & $0{,}5760$ \\
    2  & V+S+B & weighted & $0{,}8511$ & $0{,}5649$ \\
    3  & S+B   & weighted & $0{,}8779$ & $0{,}5648$ \\
    4  & V+S   & concat   & $0{,}8803$ & $0{,}5618$ \\
    5  & V+B   & concat   & $0{,}8539$ & $0{,}5613$ \\
    6  & S+B   & concat   & $0{,}8657$ & $0{,}5598$ \\
    7  & V+S   & weighted & $0{,}8695$ & $0{,}5580$ \\
    8  & V     & yok      & $0{,}8417$ & $0{,}5576$ \\
    9  & V+B   & weighted & $0{,}8582$ & $0{,}5530$ \\
    10 & V+S+B & concat   & $0{,}8520$ & $0{,}5517$ \\
    11 & B     & yok      & $0{,}8318$ & $0{,}5215$ \\
    \bottomrule
  \end{tabular}
\end{table}

\begin{table}[H]
  \centering
  \caption{0.~kat dondurulmuş $\to$ ince ayar makro-F1 farkları (ince ayar $-$
  dondurulmuş).}
  \label{tab:transfer}
  \begin{tabular}{llcc}
    \toprule
    Yapılandırma & Füzyon & $\Delta$Doğruluk & $\Delta$Makro-F1 \\
    \midrule
    Swin-T (tek)        & yok       & $+0{,}007$ & $+0{,}016$ \\
    V+S birleştirme     & concat    & $-0{,}011$ & $+0{,}011$ \\
    V+B birleştirme     & concat    & $+0{,}012$ & $+0{,}017$ \\
    S+B ağırlıklı       & weighted  & $-0{,}007$ & $+0{,}026$ \\
    V+S+B ağırlıklı     & weighted  & $+0{,}028$ & $+0{,}020$ \\
    \bottomrule
  \end{tabular}
\end{table}

İnce ayarın yakınsama hızı ve göreli maliyeti Tablo~\ref{tab:cost}'ta verilmiştir.
Her iki rejim de en iyi doğrulama makro-F1'ine birkaç epoch içinde ulaşır;
dondurulmuş rejim daha hızlıdır ve A100 gerektirmez. İnce ayarın ek maliyeti,
verimlilik ayarıyla (TF32, bfloat16, veri yükleyici paralelliği) yaklaşık
$5{,}3$--$5{,}9\times$ azaltılarak bütçe içinde tutulmuştur.

\begin{table}[H]
  \centering
  \caption{Transfer maliyeti: en iyi doğrulama epoch'u (yakınsama hızı) ve ayarlı
  A100 verim kazancı. Eğitim süresi loglanmadığından maliyet, yakınsama epoch'u ve
  göreli verimle ifade edilir.}
  \label{tab:cost}
  \resizebox{\linewidth}{!}{%
  \begin{tabular}{lccc}
    \toprule
    Yapılandırma & Dondurulmuş (en iyi epoch) & İnce ayar (en iyi epoch)
      & Ayarlı verim kazancı \\
    \midrule
    Swin-T (tek) & 4 & 6 & $\sim\!5{,}3\times$ \\
    V+S+B        & 4 & 6 & $\sim\!5{,}9\times$ \\
    \bottomrule
  \end{tabular}
  }
\end{table}

\subsection{5-katlı çapraz doğrulama}
Seçilen en iyi dört yapılandırma tam 5-katlı çapraz doğrulamada
değerlendirilmiştir (Tablo~\ref{tab:cv}, Şekil~\ref{fig:cvbar}). \textbf{Üç
omurganın ağırlıklı füzyonu} ($\text{V+S+B}$), $0{,}6094 \pm 0{,}0254$ ortalama
ve $0{,}6119$ havuzlanmış makro-F1 (\%95 GA $[0{,}5930,\,0{,}6295]$) ile en iyi
yapılandırmadır ve nihai model olarak seçilmiştir. En iyi tekli omurga
(Swin-T, havuzlanmış $0{,}5985$) ve en iyi ikili füzyon (S+B ağırlıklı,
$0{,}5893$) bunu izler. Füzyonun katkısı Bölüm~\ref{sec:discussion}'da
istatistiksel olarak da incelenmiştir.

\begin{table}[H]
  \centering
  \footnotesize
  \caption{5-katlı CV sıralaması (makro-F1 ortalamasına göre). Önyükleme \%95 GA,
  ayrık OOF test tahminleri üzerinden.}
  \label{tab:cv}
  \resizebox{\linewidth}{!}{%
  \begin{tabular}{llccccc}
    \toprule
    Yapılandırma & Füzyon & Doğruluk (ort$\pm$std) & Makro-F1 (ort$\pm$std)
      & Makro-P (ort$\pm$std) & Makro-R (ort$\pm$std) & Makro-F1 havuz [\%95 GA] \\
    \midrule
    \textbf{V+S+B} & \textbf{weighted} & $0{,}8780\pm0{,}0140$
      & $\mathbf{0{,}6094\pm0{,}0254}$ & $0{,}6111\pm0{,}0243$
      & $0{,}6246\pm0{,}0226$ & $\mathbf{0{,}6119\;[0{,}5930,\,0{,}6295]}$ \\
    Swin-T (tek)   & yok      & $0{,}8679\pm0{,}0097$ & $0{,}5969\pm0{,}0261$
      & $0{,}5958\pm0{,}0267$ & $0{,}6148\pm0{,}0257$
      & $0{,}5985\;[0{,}5806,\,0{,}6152]$ \\
    S+B            & weighted & $0{,}8652\pm0{,}0058$ & $0{,}5879\pm0{,}0111$
      & $0{,}5826\pm0{,}0138$ & $0{,}6106\pm0{,}0153$
      & $0{,}5893\;[0{,}5736,\,0{,}6054]$ \\
    V+B            & concat   & $0{,}8454\pm0{,}0084$ & $0{,}5727\pm0{,}0149$
      & $0{,}5728\pm0{,}0137$ & $0{,}6006\pm0{,}0220$
      & $0{,}5753\;[0{,}5638,\,0{,}5863]$ \\
    \bottomrule
  \end{tabular}
  }
\end{table}

\begin{figure}[H]
  \centering
  \includegraphics[width=0.6\linewidth]{cv_macrof1_bar.png}
  \caption{Seçili yapılandırmaların 5-katlı makro-F1 değerleri (ortalama $\pm$
  standart sapma). Üçlü ağırlıklı füzyon en yüksek ortalamaya sahiptir.}
  \label{fig:cvbar}
\end{figure}

\subsection{Nihai modelin ayrıntılı çözümlemesi}
Nihai modeli (fold 0) daha yakından inceliyoruz: önce eğitim dinamiği, ardından
havuzlanmış karışıklık matrisi ile sınıf bazında başarım ele alınır.

\subsubsection{Eğitim dinamiği}
Şekil~\ref{fig:curves} nihai modelin (0.~kat) eğitim eğrilerini gösterir; eğitim
kaybı düzenli azalırken doğrulama makro-F1'i erken bir epoch'ta tepe yapar ve
erken durdurma devreye girer.

\subsubsection{Karışıklık matrisi ve sınıf bazında başarım}
Şekil~\ref{fig:cm} havuzlanmış OOF karışıklık
matrisidir. Yaygın sınıflar çok yüksek köşegen değerlerine sahiptir
(ör.\ \texttt{normal-pylorus} $0{,}99$, \texttt{retroflex-stomach} $0{,}99$,
\texttt{impacted-stool} $0{,}99$), buna karşılık hataların büyük kısmı
\emph{anlamsal/anatomik olarak komşu} sınıflar arasında yoğunlaşır: ülseratif
kolit dereceleri birbirine karışır, \texttt{hemorroids} büyük ölçüde
\texttt{retroflex-rectum} olarak tahmin edilir ve \texttt{oesophagitis-a}
\texttt{normal-z-line} ile karışır.

\begin{figure}[H]
  \centering
  \begin{subfigure}{0.48\linewidth}
    \includegraphics[width=\linewidth]{training_curves.png}
    \caption{Eğitim eğrileri (0.~kat).}
    \label{fig:curves}
  \end{subfigure}\hfill
  \begin{subfigure}{0.48\linewidth}
    \includegraphics[width=\linewidth]{per_class_f1.png}
    \caption{Sınıf bazında F1 (renk = destek).}
    \label{fig:pcf1}
  \end{subfigure}
  \caption{Nihai modelin eğitim dinamikleri ve sınıf bazında başarımı. Düşük
  destekli sınıflar düşük F1 ile güçlü biçimde ilişkilidir.}
\end{figure}

\begin{figure}[H]
  \centering
  \includegraphics[width=0.92\linewidth]{confusion_matrix.png}
  \caption{Nihai modelin havuzlanmış ($n=10{.}662$) OOF karışıklık matrisi
  (satır-normalize). Hatalar büyük ölçüde anatomik/ordinal komşu sınıflar
  arasındadır.}
  \label{fig:cm}
\end{figure}

Sınıf bazında metrikler (Tablo~\ref{tab:perclass}) bu tabloyu nicelleştirir:
yüksek destekli sınıflarda F1 $\approx 0{,}97$--$0{,}99$ iken, en düşük destekli
sınıflarda (\texttt{hemorroids}: 6, \texttt{ulcerative-colitis-grade-1-2}: 11)
F1 sıfıra kadar düşer.

Havuzlanmış OOF üzerinde, ana ölçüt makro-F1'in ($0{,}612$) yanında destekleyici
toplu ölçütler de hesaplanmıştır: doğruluk $0{,}878$ (tek-etiketli çok-sınıfta
mikro-F1'e eşittir), ağırlıklı F1 $0{,}879$ ve Matthews korelasyon katsayısı
(MCC) $0{,}868$. Makro-F1 ile bu toplu ölçütler arasındaki belirgin açık,
dengesizliğin etkisini bir kez daha ortaya koyar; MCC ($0{,}868$), veri kümesi
makalesinin en iyi CNN temelinin ($0{,}899$) biraz altında olmakla birlikte aynı
bütünsel uyum bandındadır.

\begin{table}[H]
  \centering
  \caption{Nihai model: en düşük ve en yüksek F1'li altı sınıf (havuzlanmış OOF).}
  \label{tab:perclass}
  \begin{tabular}{lccc}
    \toprule
    Sınıf & Kesinlik & Duyarlılık & F1 (Destek) \\
    \midrule
    hemorroids                     & $0{,}00$ & $0{,}00$ & $0{,}00$ (6) \\
    ulcerative-colitis-grade-1-2   & $0{,}00$ & $0{,}00$ & $0{,}00$ (11) \\
    ulcerative-colitis-grade-2-3   & $0{,}05$ & $0{,}07$ & $0{,}06$ (28) \\
    \midrule
    bbps-0-1                       & $0{,}98$ & $0{,}97$ & $0{,}98$ (646) \\
    normal-pylorus                 & $0{,}98$ & $0{,}99$ & $0{,}98$ (999) \\
    retroflex-stomach              & $0{,}99$ & $0{,}99$ & $0{,}99$ (764) \\
    \bottomrule
  \end{tabular}
\end{table}

\subsection{Toplamsal teknikler ve istatistiksel anlamlılık}
\label{sec:addons}
Nihai model üzerinde denenen toplamsal, sızıntısız teknikler
Tablo~\ref{tab:addons}'da özetlenmiştir. Test zamanı veri artırma (TTA) ve
sızıntısız seed topluluğu makro-F1'i güven aralığı düzeyinde ayırt edilebilir
biçimde artırmamıştır. Son-işlem logit düzeltmesi~\cite{menon2020logitadjust}
ise makro-F1'i \emph{düşürmüştür} (Şekil~\ref{fig:logit}): düzeltme gücü $\tau$
arttıkça başarım tekdüze azalır. Bunun nedeni, modelin zaten dengeli örnekleyici
ile eğitilmiş olması ve örtük önselin yaklaşık tekdüze olmasıdır; eğitim önselini
çıkarmak \emph{ikinci kez} düzeltme yaparak nadir sınıflara aşırı kayma getirir.

Buna karşılık, eşli \textbf{McNemar} testi~\cite{dietterich1998mcnemar}
($n=10{.}662$), üçlü füzyonun en iyi tekli omurgadan ($\chi^2_{cc}=15{,}25$,
$p=9{,}2\times10^{-5}$) ve en iyi ikili füzyondan ($\chi^2_{cc}=25{,}25$,
$p=4{,}7\times10^{-7}$) \emph{anlamlı biçimde daha doğru} olduğunu göstermiştir.
Bu, doğruluk düzeyinde bir testtir; makro-F1 güven aralıklarının kısmen
örtüşmesiyle çelişmez, onu tamamlar: füzyon genel olarak anlamlı sayıda doğru
tahmin ekler, ancak makro-F1 kazancı nadir sınıflarca sınırlanır.

\begin{table}[H]
  \centering
  \caption{Nihai model üzerinde toplamsal teknikler (havuzlanmış OOF makro-F1).
  Hiçbiri nihai modeli değiştirmez; tümü sızıntısızdır.}
  \label{tab:addons}
  \footnotesize
  \begin{tabularx}{\linewidth}{@{}l c X@{}}
    \toprule
    Teknik & Havuz makro-F1 & Not \\
    \midrule
    Temel (nihai model)        & $0{,}6119$ & ana sonuç \\
    Seed topluluğu             & $0{,}6157\;[0{,}5994,\,0{,}6321]$ & küçük artış; GA örtüşüyor; final model değişmedi \\
    TTA                        & $0{,}6105\;[0{,}5920,\,0{,}6280]$ & kazanım yok; GA örtüşüyor \\
    Logit düzeltmesi ($\tau{=}1$) & $0{,}4911$ & düşürür (aşırı düzeltme) \\
    GMU füzyonu (üçlü, 0.~kat) & $0{,}5525$ & ağırlıklının altında \\
    \bottomrule
  \end{tabularx}
\end{table}

\begin{figure}[H]
  \centering
  \includegraphics[width=0.5\linewidth]{logit_adjust_tau.png}
  \caption{Son-işlem logit düzeltmesi duyarlılık eğrisi: makro-F1, düzeltme gücü
  $\tau$ ile tekdüze azalır.}
  \label{fig:logit}
\end{figure}

\subsection{Literatürle bağlamsal karşılaştırma}
Sonuçları literatürle konumlandırırken iki noktayı ayrı ele alıyoruz: çalışmalar
arasındaki protokol/ölçüt farkları ve bu farkların metrik yorumuna etkisi.

\subsubsection{Protokol ve ölçüt farkları}
Literatürdeki HiperKvasir çalışmaları farklı bölme protokolleri ve farklı ölçüt
tanımları kullandığından, sayılar doğrudan değil \emph{bağlamsal} ele alınmalıdır.
Çalışmaların çoğu doğruluk veya genel/ağırlıklı F1 raporlar; sızıntıya karşı
denetlenmiş bir makro-F1 seyrek verilir. Veri kümesini tanıtan
çalışma~\cite{borgli2020hyperkvasir} bir istisnadır: iki-bölme ortalamasıyla
\emph{makro-ortalama} F1 raporladığından ölçüt türü bizimkiyle aynıdır, fakat bölme
protokolü farklıdır. GastroViT~\cite{gastrovit2025} bizimle aynı 23 sınıflı alt
kümeyi kullanır ancak topluluk ve farklı bölme ile çalışır; EffiMix~\cite{ramamurthy2022effimix}
ve HSW-ViT~\cite{wang2023hswvit} de farklı bölmelerde sonuç verir.

\begin{table}[H]
  \centering
  \footnotesize
  \caption{HiperKvasir bağlamsal konumlandırma (protokoller farklı; doğrudan kıyas
  değil). ``$-$'': karşılaştırılabilir makro-F1 bildirilmemiş. $^{\dagger}$Borgli'nin
  doğruluğu mikro-F1 olarak verilmiştir (tek-etiketli çok-sınıfta mikro-F1, doğruluğa
  eşittir).}
  \label{tab:context}
  \begin{tabularx}{\linewidth}{@{}l X l c c l@{}}
    \toprule
    Çalışma & Yöntem & Protokol & Doğ. & Makro-F1 & Diğer \\
    \midrule
    Bu çalışma & ViT füzyonu (V+S+B) & resmi 5-katlı + GA & $\approx$0,88 & \textbf{0,612} & MCC 0,868 \\
    Borgli 2020~\cite{borgli2020hyperkvasir} & CNN (DenseNet-161) & 2-bölme ort. & 0,907$^{\dagger}$ & 0,619 & MCC 0,899 \\
    GastroViT 2025~\cite{gastrovit2025} & ViT topluluğu & kendi bölmesi & 0,920 & $-$ & $-$ \\
    HSW-ViT 2023~\cite{wang2023hswvit} & HSW-ViT & kendi bölmesi & 0,868 & $-$ & $-$ \\
    EffiMix 2022~\cite{ramamurthy2022effimix} & karışım & kendi bölmesi & 0,980 & $-$ & ağ.\ F1 0,97 \\
    \bottomrule
  \end{tabularx}
\end{table}

\subsubsection{Metrik konumlandırma}
Makro-F1 ölçütünde sonucumuz ($0{,}612$), veri kümesi makalesinin CNN temel
sonuçlarının bandındadır (makro-F1 $0{,}605$–$0{,}619$); bu, üç ViT omurgasının
füzyonunun güçlü CNN temellerine yakın bir makro başarım sağladığını bağlamsal
olarak gösterir. Yüksek doğruluk değerleri (EffiMix $\approx$\%98, GastroViT
$\approx$\%92), farklı protokollerden ve dengesizlik altında iyimser olan
toplu/ağırlıklı ölçütlerden kaynaklanır; nitekim Borgli'de de makro-F1 ($\approx
0{,}62$) ile mikro-F1/doğruluk ($\approx 0{,}91$) arasındaki açık, bu raporda
gözlenen açıkla aynı dengesizlik olgusunu yansıtır. Bu nedenle güncel-en-iyi bir
başarı iddiasında bulunulmamış; konum bağlamsal olarak sunulmuştur.
